% BC-G

BC-G implements dynamic load-balancing by using the GLB library. Note
that this is a new capability (over what is supported by BC) -- therefore it is not surprising that
new code needs to be written and the line count is actually higher
than for BC. Specifically, \code{BCTaskFrame} and \code{BCTaskBag}
needs to be implemented, together with the invocation of the
\code{GLB} framework. A task captures a (start, end) range of
vertices. The code for \code{BCTaskFrame} is
straightforward -- it pops tasks from the \code{BCTaskBag} and for
each task executes it by calling \code{bfsShortestPaths} already
implemented in BC. The reducer takes two betweenness maps (vectors) and sums
them up pointwise.

We now consider the implementation of \code{BCTaskBag}. Our current
implementation (reported in this submission) is preliminary and in
need of further work. It is designed to satisfy the following
properties. It should be possible to find the ``mid point'' fast so
that the bag can be split in two quickly. The cost of a merge should
be low. 

The current implementation uses an array, a head pointer and a tail
pointer. The head pointer is moved forward whenever the task frame
retrieves an item and works on it. The tail pointer is used during
splitting, so it can move to the split point, copy the latter half of
the data and send it out. 

When merging, we create a new array and copy the local taskbag and the
incoming taskbag to the new array and set up the head and tail pointer
properly. We believe this copy is expensive and should be eliminated.

We believe a more careful design for this data-structure will improve
performance. Basically we need to design a data-structure that
represents a set of intervals such that two such sets can be combined
quickly and split quickly, where the split ensures that the sum of the
sizes of intervals in each half are roughly equal.  We believe that
improving the design of this data-structure will significantly improve
performance. {\em We expect to work on this in the next few
  weeks, before SuperComputing.}

\subsection{Peformance measurement}

Section~\ref{sec:bc} presents the background for the Betweenness
Centrality benchmark. There it was pointed out that the static
partitioning of the work across the given number of places may not be
ideal.  

As described in the benchmark, the application metric traversed edges
per second (TEPS) is measured as:

$$ TEPS = (8\times N \times N)/t$$

\noindent where $N$ is the number of vertices.

We now discuss results. 
All tests were run with particular parameters $(1=0.8, b=c=0.05,
d=0.1)$. This gives rise to a significantly unbalanced graph.

For Power~775 we had access to a small number of cores, and the machine was
in heavy demand, so we were unable to run long jobs.  So we chose to
vary $n$ (the log of the number of vertices) as we varied $p$. The
data is presented in a table form in Table~\ref{table:perf:bc-p7ih} to
show that $n$ varies as $P$ varies. The data shows that at these
scales, and on this architecture BC-G is quite effective in reducing
the runtime (for a given $P$ and $n$), e.g. achieving a $2x$ speed-up
for BC-G over BC.

For Blue Gene/Q and K, we show charts for a fixed $n=19$ 
(Figure~\ref{fig:perf:bc-bgq} and Figure~\ref{fig:perf:bc-k}). 
 For Blue
Gene/Q the data shows that BC-G is able to achieve around a 2x speedup
for most values of $P$. For $P=2048$, however, the performance is
somewhat comparable. For K, BC-G is generally better than BC, but with
a smaller constant. 

We believe that the representation of \code{BCTaskBag} can be
improved. The \code{split} and \code{merge} operations are expensive
and this may account for why the speedup worsens with larger $P$
because there is more contention and stealing.  {\em We expect to work
  on improving these numbers for the presentation at SC, if selected.}


\begin{figure}
\centering

$$
\begin{array}{|l|l|l|l|}
\hline
\mbox{P}&\mbox{n}&\mbox{BC (s)}&\mbox{BC-G (s)}\\ \hline \hline
1 & 16 & 78 & 109 \\ \hline
2 & 16 & 32 & 34  \\ \hline
4 & 16 & 21 & 17 \\ \hline
8 & 16 & 13 & 9 \\ \hline
16 & 17 & 35 & 19 \\ \hline
32 & 17 & 21 & 11 \\ \hline 
64 & 18 & 67 & 34 \\ \hline 
128 & 18 & 41 & 20 \\ \hline 
256 & 19 & 70 & 55 \\ \hline
\end{array}
$$
\caption{X10 Performance for BC (on Power~775)}
\label{table:perf:bc-p7ih}

\includegraphics[width=3.25in]{figures/bgq/bc-bgq.eps}
\caption{X10 Performance for BC (on Blue~Gene/Q)}
\label{fig:perf:bc-bgq}

\includegraphics[width=3.25in]{figures/k/bc-k.eps}
\caption{X10 Performance for BC (on K)}
\label{fig:perf:bc-k}
\end{figure}




